% p2-08-a5-flavor

\section{P2 §8. The A5 Flavor Symmetry Boundary}

8. The A5 Flavor Symmetry Boundary

         Figure (fig-p2-ch8-a5-boundary). A5-boundary triptych --- Temptation / Missing morphism /
                                             Boundary marker.

  8.1 The Temptation
  The icosahedral group A5 cong PSL(2,5) has order 60 and is isomorphic to H3 modulo
  sign. Because A5 involves 5 in its character table and $\varphi$ = (1+5)/2, there is a superficial
  temptation to connect A5 flavor symmetry models to the $\varphi$-structured patterns seen in
  Paper 1. The A5 group has been studied as a candidate discrete flavor symmetry for
  explaining golden-ratio neutrino mixing angles (Ding, Everett, and Stuart, Nucl. Phys. B
  854 (2012) 291, ScienceDirect), where the solar mixing angle is predicted by
  cottheta12 = $\varphi$.

  8.2 Why This Connection Is Not Made Here [L4 --- Candidate; but
  distinct from Paper 1 subject matter]

  The flavor symmetry group acting on lepton doublets in discrete flavor models is a
  discrete subgroup of the continuous symmetry that would rotate the three generations.
  The golden-ratio mixing scheme predicts cottheta12 = $\varphi$, with the atmospheric mixing
  angle satisfying tantheta23 = 1 and the reactor angle constrained by vacuum
  alignment. This is a different physical context from the coupling constant
  compressibility of Paper 1, and no formal connection between the two has been
  established.

  Explicit caveats:

  1. [L4 $\rightarrow$ watch L6] A5 flavor symmetry models require a vacuum alignment mechanism;
  the breaking scale is not fixed by the symmetry alone. Daya Bay and RENO
  measurements of theta13 approx 8.$9^{\circ}$ strongly constrain golden-ratio mixing variants;
  models must accommodate this with next-to-leading-order corrections.

  2. The SU(3) color group of QCD has no structural connection to A5 or $\varphi$. Color charge
  is not $\varphi$-structured in any derivable sense. No claim of any kind is made about a QCD-$\varphi
  connection in this paper.

  3. Even if A5 flavor symmetry were confirmed, the connection to the E8/Toda $\varphi
  occurrence would require an explicit embedding construction linking the neutrino
  sector symmetry-breaking to the mass-spectrum structure of 2D integrable theories. No
  such construction exists.

  4. The formal success/failure criterion for the A5 mechanism is specified in Section 4.4.

\subsection{References}

- Coldea, R. et al. "Quantum Criticality in an Ising Chain: Experimental Evidence for Emergent E8 Symmetry." \textit{Science} 327 (2010) 177. DOI \href{https://doi.org/10.1126/science.1180085}{10.1126/science.1180085}.
- Wu, J. et al. "E8 Symmetry in CoN$b_{2}$$O_{6}$ at Higher Resolution." \textit{Phys. Rev. B} 108, L060403 (2023). DOI \href{https://doi.org/10.1103/PhysRevB.108.L060403}{10.1103/PhysRevB.108.L060403}.
- Fring, A.; Korff, C. "Affine Toda field theories related to Coxeter groups of non-crystallographic type." \textit{Nucl. Phys. B} 729 (2005) 361. \href{https://arxiv.org/abs/hep-th/0506226}{arXiv:hep-th/0506226}. DOI \href{https://doi.org/10.1016/j.nuclphysb.2005.08.044}{10.1016/j.nuclphysb.2005.08.044}.
- Pellis, S. "The Fine-Structure Constant and the Golden Ratio." \href{https://vixra.org/pdf/2110.0117v4.pdf}{viXra 2110.0117 v4}.
